\section{Related Work}

\textbf{Deobfuscation and robustness.}
Prior work on code deobfuscation focuses primarily on recovering a cleaner
representation from transformed code. Our setting instead evaluates program
reasoning while the code remains obfuscated, measuring robustness to
semantics-preserving changes rather than reconstruction quality.

\textbf{Model merging.}
Model merging can preserve or strengthen knowledge shared across independently
trained models, while knowledge specific to individual models is more likely to
degrade during composition~\cite{zaman-etal-2024-fuse,hess2025forgetting}.
This asymmetry has also been exploited for bias reduction and machine
unlearning~\cite{zaman-etal-2024-fuse,xu-etal-2025-zjuklab}.
Merge quality further depends on specialist training dynamics: memorization late
in fine-tuning can produce idiosyncratic parameter updates that interfere during
merging, while earlier checkpoints can compose more effectively
~\cite{horoi2026memorization}.

\textbf{Modular adaptation.}
Related work identifies similar shared-versus-specific structure within
low-rank adapters. LoRA can exhibit capacity limits on high-entropy objectives,
motivating architectures that route updates between low-rank and full-parameter
experts~\cite{tang2026beyond}. Token-dependent routing has likewise been used to
separate shared and subdomain-specific adapter updates
~\cite{chen2026alter}.

\textbf{Connection to our setting.}
These findings motivate our treatment of obfuscation specialists as containing
both transform-specific adaptations and reasoning updates shared across
semantics-preserving views. We test whether parameter-space composition can
retain this shared program-reasoning capability while reducing dependence on
individual transformation surfaces.